%------------------------------------------------------------------------------%
\documentclass[a4paper,12pt,fleqn]{article}

\usepackage{calculo_varias_variaveis-1}

\ShowAnswers

%------------------------------------------------------------------------------%
\begin{document}

\makeHeader{CFVVI}{Prova 1 -- Turma 7}{05/05/2025}

% 01
%------------------------------------------------------------------------------%
\exercise{25}
Encontre uma parametrização para o movimento de uma partícula que parte do
ponto \((a, 0)\) e traça o círculo \(x^2+y^2=a^2\) duas vezes no sentido horário.
\clearpagequestiononly

\begin{answer}
  Duas voltas no círculo de raio $a$ centrado na origem

  Parametrização do círculo unitário centrado na origem no sentido anti-horário
  \begin{align*}
    x      & = \cos(\theta) \\
    y      & = \sin(\theta) \\
    \theta & \in [0, 2\pi]
  \end{align*}
  Parametrização do círculo de raio $a$ centrado na origem
  \begin{align*}
    x      & = a\cos(\theta) \\
    y      & = a\sin(\theta) \\
    \theta & \in [0, 2\pi]
  \end{align*}
  Realizando duas voltas no círculo
  \begin{align*}
    x      & = a\cos(\theta) \\
    y      & = a\sin(\theta) \\
    \theta & \in [0, 4\pi]
  \end{align*}
  Invertendo o sentido de rotação
  \begin{align*}
    x      & = a\cos(4\pi -\theta) = a\cos(-\theta) = \hpm a\cos(\hpm\theta) \\
    y      & = a\sin(4\pi -\theta) = a\sin(-\theta) = -    a\sin(   -\theta) \\
    \theta & \in [0, 4\pi]
  \end{align*}
\end{answer}

% 02
%------------------------------------------------------------------------------%
\exercise{25}
Considerando a função
\(
  f(x,y) = \sqrt{4 - x - y^2}
\)
\begin{enumerate}[label=(\alph*)]
  \item \ [10] Determine o domínio de $f$ e represente-o graficamente
  \item \ [5]  Encontre a imagem de $f$
  \item \ [10] Caracterize todas as curvas de nível de $f$ e esboce três delas
\end{enumerate}
\begin{questiononly}
  \begin{multicols}{2}
  Domínio
  \vspace{-0.8\baselineskip}
  \begin{center}
    \includegraphics{axis_xy}
  \end{center}
  Curvas de nível
  \vspace{-0.8\baselineskip}
  \begin{center}
    \includegraphics{axis_xy}
  \end{center}
  \end{multicols}
  \clearpage
\end{questiononly}

\begin{answer}
\begin{multicols}{2}
  Domínio
  \vspace{-0.8\baselineskip}
  \begin{center}
    \includegraphics{P1_T7_dominio}
  \end{center}
  Curvas de nível
  \vspace{-0.8\baselineskip}
  \begin{center}
    \includegraphics{P1_T7_curvas_nivel}
  \end{center}
\end{multicols}

\begin{enumerate}[label=(\alph*)]
\item O domínio de $f$ é o conjunto de pontos \((x, y)\) tais que
  \begin{gather*}
    4 - x - y^2 \geq 0 \\
    x \leq 4 - y^2
  \end{gather*}
  Portanto, o domínio é
  \[
    D_f = \{(x, y) \in \R^2 \st x \leq 4 - y^2\}
  \]
  Geometricamente, corresponde à região à esquerda da parábola \( x = 4 - y^2 \)

\item Como \( 4 - x - y^2 \in [0, 4] \) no domínio, temos
  \(
    \text{Im}(f) = [0, 2]
  \)

\item As curvas de nível são determinadas por \( f(x,y) = c \) para \( c \in [0, 2] \):
  \begin{gather*}
    \sqrt{4 - x - y^2} = c \\
    x = 4 - y^2 - c^2
  \end{gather*}
  isto é, parábolas voltadas para a esquerda
  \begin{align*}
    c & = 0 & x & = 4 - y^2 -0^2 = 4 - y^2 \\
    c & = 1 & x & = 4 - y^2 -1^2 = 3 - y^2 \\
    c & = 2 & x & = 4 - y^2 -2^2 =   - y^2
  \end{align*}
\end{enumerate}
\end{answer}

% 03
%------------------------------------------------------------------------------%
\exercise{30}
Calcule o limite ou mostre que ele não existe
\begin{enumerate}[label=(\alph*)]
  \item \ [15] \(\lim_{(x,y) \to (0,0)} \frac{xy^2}{x^2 + y^4}\)
  \item \ [15] \(\lim_{(x,y) \to (0,1)} \frac{\ln(y)\sin(x)}{x}\)
\end{enumerate}
\clearpagequestiononly

\begin{answer}
  \begin{enumerate}[label=(\alph*)]
    \item Testando diferentes caminhos \\
    Caminho \(x = y^2\)
    \[
    \lim_{(x,y) \to (0,0)} \frac{xy^2}{x^2 + y^4} \evalat{x = y^2}{}
    = \lim_{y \to 0} \frac{y^2 y^2}{y^4 + y^4}
    = \lim_{y \to 0} \frac{y^4}{2y^4}
    = \lim_{y \to 0} \frac{1}{2}
    = \frac{1}{2}
    \]
    Caminho \(x = 0\)
    \[
    \lim_{(x,y) \to (0,0)} \frac{xy^2}{x^2 + y^4} \evalat{x = 0}{}
    = \lim_{y \to 0} \frac{0 \times y^2}{0 + y^4}
    = \lim_{y \to 0} 0
    = 0
    \]
    Como os limites são diferentes, o limite \textbf{não existe}
    \item
    \begin{align*}
      \lim_{(x,y) \to (0,1)} \frac{\ln(y)\sin(x)}{x}
      & = \left(\lim_{(x,y) \to (0,1)} \ln(y)\right)
          \left(\lim_{(x,y) \to (0,1)} \frac{\sin(x)}{x}\right) \\
      & = \ln(1)
          \times 1 \\
      & = 0
    \end{align*}
  \end{enumerate}
\end{answer}

% 04
%------------------------------------------------------------------------------%
\exercise{20}
Encontre e esboce as curvas que representam o corte da superfície
\(4x^2+y^2=16\)
pelos planos
\(z=0\) e
\(y=0\)
\begin{questiononly}
  \begin{multicols}{2}
    \(z=0\)
    \vspace{-0.8\baselineskip}
    \begin{center}
      \includegraphics{axis_xy}
    \end{center}
    \(y=0\)
    \vspace{-0.8\baselineskip}
    \begin{center}
      \includegraphics{axis_xz}
    \end{center}
  \end{multicols}
\end{questiononly}

\begin{answer}
  \begin{multicols}{2}
    \(z=0\)
    \vspace{-0.8\baselineskip}
    \begin{center}
      \includegraphics{T7_conica_z0}
    \end{center}
    \(y=0\)
    \vspace{-0.8\baselineskip}
    \begin{center}
      \includegraphics{T7_conica_y0}
    \end{center}
  \end{multicols}

  No plano \(z=0\) a equação \(4x^2+y^2=16\) permanece inalterada, porém,
  agora representa a elipse
  \[
    \frac{x^2}{2^2} + \frac{y^2}{4^2} = 1
  \]
  que intercepta o eixo $x$ nos pontos \((2, 0)\) e \((-2, 0)\)
  e o eixo $y$ nos pontos \((4, 0)\) e \((-4, 0)\)

  No plano \(y=0\) a equação \(4x^2+y^2=16\) se reduz a
  \begin{align*}
    4x^2 + 0 & = 16    \\
    x^2      & = 4     \\
    \abs{x}  & = 2     \\
    x        & = \pm 2
  \end{align*}
  que representa as retas verticais \(x=2\) e \(x=-2\)
\end{answer}

%------------------------------------------------------------------------------%
\end{document}
%------------------------------------------------------------------------------%
